\chapter{Разбиение выборки, метрики и базовые модели}

\chaptergoals
  {Объяснить, почему оценка модели требует независимой схемы проверки, зачем нужен baseline и как метрики качества переводятся в содержательный вывод для регрессионной и временной постановки.}
  {
    \item разделять данные на обучающую и тестовую части в соответствии с задачей;
    \item объяснять назначение baseline-модели;
    \item рассчитывать и интерпретировать метрики \(MAE\), \(RMSE\) и \(R^2\);
    \item различать обычное случайное разбиение и хронологическую схему проверки для временных рядов.
  }

\section{Почему нельзя оценивать модель на тех же данных, на которых она обучалась}

Если качество модели рассчитывается на той же выборке, по которой она подбирала параметры, полученная оценка почти всегда будет слишком оптимистичной. В таком случае модель проверяется не на способности обобщать, а на способности воспроизводить уже увиденные наблюдения. Для инженерной практики это неприемлемо, поскольку интерес представляет поведение модели на новых данных.

Поэтому в анализе данных вводится разбиение выборки. В простейшем случае используются обучающая и тестовая части. Обучающая часть нужна для настройки модели, тестовая --- для независимой проверки ее качества. В более строгих постановках добавляется валидационная выборка, но для первых учебных примеров достаточно базовой схемы сравнения.

\section{Зачем нужен baseline}

До выбора более содержательной модели необходимо ввести baseline --- простейший ориентир, который задает минимально допустимый уровень качества. Его роль методически принципиальна. Без такого ориентира можно получить численное значение метрики и не понять, является ли оно убедительным.

Для регрессии естественным baseline-решением является прогноз средним значением по обучающей выборке:
\[
\hat{y}_i = \overline{y}_{train}.
\]

Смысл этой конструкции прост. Если более сложная модель не улучшает такой примитивный прогноз, значит, либо признаки не содержат полезной информации, либо сама модель выбрана неудачно. В обоих случаях усложнение без дополнительного анализа не имеет исследовательского смысла.

\section{Метрики регрессии и их интерпретация}

В настоящем пособии для первых регрессионных примеров используются три базовые метрики.

\subsection*{Средняя абсолютная ошибка}

\[
MAE = \frac{1}{n}\sum_{i=1}^{n} |y_i - \hat{y}_i|.
\]

Эта метрика показывает среднее абсолютное отклонение прогноза от фактического значения. Ее преимущество состоит в том, что она измеряется в тех же единицах, что и прогнозируемая величина. Поэтому \(MAE\) удобно обсуждать в инженерной интерпретации: можно сразу сказать, насколько велика типичная ошибка в мегаваттах, киловаттах или другой физической единице.

\subsection*{Корень из средней квадратичной ошибки}

\[
RMSE = \sqrt{\frac{1}{n}\sum_{i=1}^{n}(y_i - \hat{y}_i)^2}.
\]

По сравнению с \(MAE\) эта метрика сильнее реагирует на крупные ошибки. Если в задаче особенно нежелательны единичные большие промахи прогноза, \(RMSE\) дает более чувствительную оценку риска.

\subsection*{Коэффициент детерминации}

\[
R^2 = 1 - \frac{\sum_{i=1}^{n}(y_i - \hat{y}_i)^2}{\sum_{i=1}^{n}(y_i - \overline{y})^2}.
\]

Коэффициент \(R^2\) показывает, какую долю вариации целевой переменной объясняет модель относительно прогноза средним значением. Эта метрика удобна для общего сравнения, но ее не следует трактовать отдельно от абсолютных ошибок. В инженерной задаче модель может иметь приемлемый \(R^2\), но при этом допускать слишком большую ошибку в физических единицах.

\section{Как связывать метрики с содержательным выводом}

Метрика полезна лишь тогда, когда она переведена на язык задачи. Если \(MAE\) мало, нужно пояснить, мало ли оно в абсолютном масштабе установки или режима. Если \(RMSE\) заметно превышает \(MAE\), это может указывать на наличие отдельных крупных промахов. Если \(R^2\) невелик, следует обсудить, достаточно ли признаков для описания процесса.

Именно поэтому в практикуме после таблицы метрик строится диаграмма ``факт -- прогноз''. Она позволяет увидеть не только средний уровень ошибки, но и форму отклонений: систематически ли модель завышает прогноз, есть ли диапазоны, где она ошибается сильнее, и насколько плотно точки группируются вдоль диагонали идеального совпадения.

\screenshotplaceholder
  {fig:baseline-metrics}
  {Сравнение baseline-модели и линейной регрессии}
  {Таблица метрик из блокнота \code{04\_splitting\_metrics\_and\_baselines.ipynb}, где сравниваются baseline и линейная регрессия по \(MAE\), \(RMSE\) и \(R^2\).}
  {Выполнить блок сравнения моделей в \code{04\_splitting\_metrics\_and\_baselines.ipynb} и сделать скриншот итоговой таблицы так, чтобы были различимы строки baseline и линейной регрессии, а также значения метрик.}

\figureinstruction
  {Сравнение качества baseline-модели и линейной регрессии на учебном примере.}
  {После рисунка следует подчеркнуть, что baseline нужен не для ``хорошего'' качества сам по себе, а как обязательная точка отсчета. Только при наличии такого ориентира можно аргументированно утверждать, что модель действительно извлекает полезную информацию из признаков.}

\section{Особенности временных рядов}

Для временных рядов случайное разбиение выборки обычно неприемлемо. Оно разрушает естественный порядок наблюдений и может привести к утечке будущей информации в обучение. Поэтому в задачах прогноза необходимо использовать хронологическую схему: модель обучается на прошлом, а проверяется на более позднем интервале.

Из этого следует важный методический вывод. Одна и та же техника оценки качества не переносится автоматически между задачами. В регрессии по независимым наблюдениям допустимо случайное разбиение, а во временных рядах требуется уважать структуру времени.

\section{Типичные ошибки на этапе оценки качества}

Наиболее распространенные ошибки на этом этапе таковы:
\begin{itemize}[leftmargin=1.25cm]
\item отсутствие baseline и, как следствие, отсутствие точки сравнения;
\item преждевременное обращение к тестовой выборке при подборе модели;
\item использование только одной метрики без содержательной интерпретации;
\item перенос случайной схемы разбиения на временной ряд;
\item попытка обсуждать качество модели без связи с физическим масштабом ошибки.
\end{itemize}

Избежать этих ошибок помогает простое правило: метрики и схема проверки должны быть согласованы с природой задачи, а итоговый вывод должен объяснять не только число, но и его инженерный смысл.

\section{Связь главы с практикумом}

Главе соответствует практикум по сравнению baseline и линейной модели. В репозитории он доступен как \notebookhref{04_splitting_metrics_and_baselines.ipynb}{\code{04\_splitting\_metrics\_and\_baselines.ipynb}}. На реальном наборе данных по мощности электростанции в нем последовательно демонстрируются постановка регрессионной задачи, разбиение выборки, baseline, линейная модель, таблица метрик и диаграмма сравнения фактических и прогнозных значений. Благодаря этому глава остается теоретически цельной, но не отрывается от исполняемого практикума.
